% theory_environments.tex —— MIPSolvers 模块手册"通用理论层"共享环境
% 用途：为 docs/modules/ 各手册的 theory_*.tex 理论章节提供统一排版与标记机制。
% 移植自 HybridACDCDistributionSystemsSimulation/docs/_manual_common/theory_environments.tex。
% 引入方式（二选一，不得重复引入）：
%   (a) 使用公共样式的手册：mipsolvers_manual.sty 已自动引入，无需任何操作；
%   (b) 自带导言区的手册：在导言区写
%       \input{../../_manual_common/theory_environments.tex}
% 规范见 docs/modules/THEORY_WRITING_GUIDE.md 与 docs/modules/README.md。

\RequirePackage{amsthm}
\RequirePackage[most]{tcolorbox}

% ── 定理族（中文名，按 section 编号）────────────────────────────
\theoremstyle{plain}
\newtheorem{theorem}{定理}[section]
\newtheorem{lemma}[theorem]{引理}
\newtheorem{proposition}[theorem]{命题}
\newtheorem{corollary}[theorem]{推论}
\theoremstyle{definition}
\newtheorem{definition}[theorem]{定义}
\newtheorem{example}[theorem]{例}
\theoremstyle{remark}
\newtheorem{remark}[theorem]{注记}
\renewcommand{\proofname}{证明}

% ── 理论章声明框：每个 theory_*.tex 章首必须使用 ─────────────────
% 用法：\begin{theorynote} 本章为通用数学理论……\end{theorynote}
\newtcolorbox{theorynote}{%
  enhanced, breakable, colback=blue!4, colframe=blue!45!black,
  boxrule=0.6pt, arc=1.5mm, left=2mm, right=2mm,
  title={\textbf{通用理论章说明}},
  fonttitle=\small\bfseries}

% ── 未实现/理论扩展标记框：理论内容超出当前代码实现时必须就地标记 ──
% 用法：\begin{gapnote} 本节模型在代码中尚未实现……\end{gapnote}
\newtcolorbox{gapnote}{%
  enhanced, breakable, colback=orange!6, colframe=orange!70!black,
  boxrule=0.6pt, arc=1.5mm, left=2mm, right=2mm,
  title={\textbf{未实现扩展说明}},
  fonttitle=\small\bfseries}

% ── 模型边界说明：用于明确拒绝或适用范围，不表示待实现计划 ──────
\newtcolorbox{boundarynote}{%
  enhanced, breakable, colback=gray!5, colframe=gray!65!black,
  boxrule=0.6pt, arc=1.5mm, left=2mm, right=2mm,
  title={\textbf{模型边界说明}},
  fonttitle=\small\bfseries}

% ── 实现状态标签（"与实现的对应"表专用）─────────────────────────
% \implfull{文件:函数名}  已实现（附锚点，禁止行号）
% \implpart{差异说明}     部分实现（必须说明差异）
% \implnone               未实现
\newcommand{\implfull}[1]{\textbf{已实现}\quad\srcpath{#1}}
\newcommand{\implpart}[1]{\textbf{部分实现}\quad（#1）}
\newcommand{\implnone}{\textbf{未实现}}
